\scanpage{38}\mspage{31}
% archive sheet number: 37
comme un \uline{adjoint à droite} (pl. fidèle) du foncteur de localisation
\[
 \pi_{d}M_{c}\longrightarrow(\pi_{d}M_{c})W_{c}^{-1}\approx\operatorname{Ho}_{W},
\]
et $\operatorname{Ho}_{W}$ se déduit de $\pi_{d}M_{c}$ par un \uline{calcul de fractions} \uline{à gauche}. Dualement, si on a exactitude à droite, \deleted{alors} i.e.
\[
 TF=F\cap W\quad(\Longleftrightarrow\quad TF=F\cap W^{\deleted{\unsure{\mathrm{univ}|\mathrm{uni}}}}_{\mathrm{univ}}),
\]
alors la relation d'équivalence pertinente sur $\operatorname{Hom}_{M}(X,Y)$, pour $X,Y$ dans $M_{cf}$, est la Q-\uline{homotopie à gauche} (impliquée par la Q-homotopie à droite), et cette fois la flèche \unsure{naturelle} de $\pi_{g}M_{cf}=\pi_{\deleted{\unsure{g}}}M_{cf}\approx\operatorname{Ho}_{W}$ dans $\pi_{d}M_{f}$
\diagram{\begin{tikzcd}[ampersand replacement=\&,row sep=small]
\pi M_{cf}\arrow[r] \& \pi_{d}M_f\\
\operatorname{Ho}_{W}\arrow[u,phantom,"\approx"{description,sloped}] \&
\end{tikzcd}}
est un \uline{adjoint à gauche} \inserted{(pl. fidèle)} du foncteur de localisation
\[
 \pi_{\unsure{d}}M_{f}\longrightarrow M_{f}W_{f}^{-1}\approx\operatorname{Ho}_{W},
\]
et \deleted{\unsure{la}}\unsure{sa|la} \unsure{localisée|localisation} s'explicite par un calcul de fractions à droite. \deleted{Dans le}

Dans le cas biexact (i.e. \inserted{celui} des véritables structures de Quillen), les deux relations de Q-homotopie dans $M_{cf}$ (\deleted{et} plus génér. dans \unsure{un} $\operatorname{Hom}_{M}(X,Y)$, si $X$ dans $M_{c}$ et $Y$ dans $M_{f}$)
